\chapter{Crawler e extração de informações}

O crawler foi construído em Scrapy com extratores especializados. O objetivo não é simplesmente copiar texto, mas produzir observações estruturadas, auditáveis e vinculadas à página exata.

\section{Configurações responsáveis}

\begin{lstlisting}
ROBOTSTXT_OBEY = True
AUTOTHROTTLE_ENABLED = True
CONCURRENT_REQUESTS_PER_DOMAIN = 2
DOWNLOAD_DELAY = 1.5
DEPTH_LIMIT = 2
USER_AGENT = "HEPConferenceHubBot/0.1 (+contato-do-projeto)"
\end{lstlisting}

\section{Descoberta semântica de páginas}

O módulo \texttt{page\_discovery.py} classifica links por função. Exemplos: chamada de resumos, datas importantes, inscrição, autores, programa, proceedings e newsletter. A navegação permanece restrita ao domínio oficial ou ao prefixo confiável do evento em uma plataforma Indico.

\section{Extração de prazos}

O extrator identifica candidatos rotulados, preservando tipo, data, horário, fuso, estado, URL, trecho e confiança. Datas sem rótulo suficiente são ignoradas para evitar falsos positivos.

\section{Conflitos}

Quando páginas oficiais divergem, os valores são mantidos como evidências conflitantes. O sistema não escolhe silenciosamente uma data. Esse comportamento foi usado em casos como diferenças entre “Important Dates” e formulários Indico.

\section{Spiders}

\begin{center}
\small
\renewcommand{\arraystretch}{1.15}
\begin{longtable}{>{\RaggedRight\arraybackslash}p{0.29\textwidth}>{\RaggedRight\arraybackslash}p{0.64\textwidth}}
\toprule
\textbf{\textbf{Spider}} & \textbf{\textbf{Finalidade}} \\
\midrule
\endfirsthead
\toprule
\textbf{\textbf{Spider}} & \textbf{\textbf{Finalidade}} \\
\midrule
\endhead
\midrule
\multicolumn{2}{r}{\footnotesize Continua na próxima página} \\
\endfoot
\bottomrule
\endlastfoot
ichep\_2026 & Coleta específica do ICHEP e páginas internas. \\
event\_site & Coletor genérico parametrizado por URL e identificador. \\
latin\_america\_events & Percorre o catálogo regional verificado. \\
latin\_america\_discovery & Descobre candidatos em fontes regionais, sem publicação automática. \\
\end{longtable}
\end{center}

\section{Saídas}

\begin{lstlisting}
output/pages.jsonl
output/events_aggregated.json
output/latin_america_pages.jsonl
output/latin_america/events_aggregated.json
output/latin_america_candidates.jsonl
\end{lstlisting}
